\section{Method}

\subsection{Overview}

DecayLinearAttention replaces standard attention with a **chunked, 
data-dependent linear recurrence**. For each head, we learn a per-token decay 
$\lambda_{h,t}$ and propagate a running state:

\[
S_t = \lambda_{h,t} S_{t-1} + \phi(k_t) v_t^\top
\tag{4.1}
\]

\[
z_t = \lambda_{h,t} z_{t-1} + \phi(k_t)
\tag{4.2}
\]

The output:

\[
O_t = \frac{\phi(q_t)^\top S_t}{\phi(q_t)^\top z_t + \epsilon}
\tag{4.3}
\]

\textbf{Figure 4.1} shows the pipeline.

\begin{lstlisting}
    X $\in$ R^{B$\times$L$\times$d}
         |
    +----+----+----+
                    
    W_q W_k W_v  decay_gate
    |   |   |    |
                  
    Q   K   V   $\lambda$_t
    |   |   |    |
                 |
     (Q)  (K) V  |
    |   |   |    |
    +---+---+----+
             
    Chunked Scan (S, z)
        |
         
    O =  (Q) S / ( (Q) z)
        |
         
    W_o
        |
         
    Output Y
\end{lstlisting}

\subsection{Data-Dependent Decay}

For each head $h$ and token $t$:

\[
\tilde{\lambda}_{h,t} = W_g x_t + b_g
\tag{4.4}
\]

\[
\lambda_{h,t} = \sigma\!\left(\tilde{\lambda}_{h,t} + \ell_h\right)
\tag{4.5}
\]

where $\ell_h$ is a \textbf{learned scalar} per head.

\textbf{Initialization.}
\begin{itemize}
\item $W_g = 0$ (no input dependence initially).
\item $b_g = 0$.
\item $\ell_h$ is initialized to linearly spaced values in
  $[{{LOG_DECAY_MIN}}, {{LOG_DECAY_MAX}}]$:

\end{itemize}
\begin{lstlisting}
  self.log_decay = nn.Parameter(torch.linspace(-0.5, 0.0, num_heads))
\end{lstlisting}

\textbf{Result.} At step 0, different heads start with different horizons:
\begin{itemize}
\item Head 0: $\lambda \approx \sigma({{LOG_DECAY_MIN}}) \approx {{LAMBDA_MIN}}$
\item Head $H-1$: $\lambda \approx \sigma({{LOG_DECAY_MAX}}) \approx {{LAMBDA_MAX}}$

\end{itemize}
This head diversity is analogous to multi-scale temporal receptive fields in 
convolutional networks.

\subsection{ELU+1 Feature Map}

\[
\phi(x) = \text{elu}(x) + 1
\tag{4.6}
\]

Applied element-wise to Q, K.

\textbf{Why ELU+1?}
\begin{itemize}
\item \textbf{Positive:} Ensures $z_t > 0$.
\item \textbf{Cheap:} No exponential or trigonometric functions.
\item \textbf{Established:} Used by Linear Transformer (Katharopoulos et al., 2020).

\end{itemize}
\subsection{Chunked Scan}

For a chunk of size $C$ starting at position $s$:

\subsubsection{Cumulative Log-Decay}

Within the chunk, we compute:

\[
\Lambda_i = \sum_{p=s}^{s+i} \log \lambda_{h,p}, \quad i \in [0, C)
\tag{4.7}
\]

via \texttt{torch.cumsum}.

\subsubsection{Intra-Chunk Attention}

For positions $i, j$ in the chunk with $j \leq i$:

\[
D_{i, j}^\text{intra} = \exp(\Lambda_i - \Lambda_j)
\tag{4.8}
\]

\[
A_{i, j}^\text{intra} = \phi(q_i)^\top \phi(k_j) \cdot D_{i, j}^\text{intra}
\tag{4.9}
\]

with $A_{i, j}^\text{intra} = 0$ for $j > i$ (causality).

Output (numerator):

\[
\text{num}_i^\text{intra} = \sum_{j \leq i} A_{i, j}^\text{intra} \cdot v_j
\tag{4.10}
\]

Normalizer (denominator):

\[
\text{denom}_i^\text{intra} = \sum_{j \leq i} A_{i, j}^\text{intra}
\tag{4.11}
\]

\subsubsection{Inter-Chunk Contribution}

Let $S_\text{prev}$, $z_\text{prev}$ be the state before the chunk. Define:

\[
D_i^\text{inter} = \exp(\Lambda_i)
\tag{4.12}
\]

Then:

\[
\text{num}_i^\text{inter} = \phi(q_i)^\top S_\text{prev} \cdot D_i^\text{inter}
\tag{4.13}
\]

\[
\text{denom}_i^\text{inter} = \phi(q_i)^\top z_\text{prev} \cdot D_i^\text{inter}
\tag{4.14}
\]

\subsubsection{Output}

\[
O_i = \frac{\text{num}_i^\text{intra} + \text{num}_i^\text{inter}}{\text{denom}_i^\text{intra} + \text{denom}_i^\text{inter} + \epsilon}
\tag{4.15}
\]

\subsubsection{State Update}

After the chunk, the state is updated for the next chunk:

\[
S_\text{new} = \exp(\Lambda_C) S_\text{prev} + \sum_{i=0}^{C-1} \exp(\Lambda_C - \Lambda_i) \phi(k_i) v_i^\top
\tag{4.16}
\]

\[
z_\text{new} = \exp(\Lambda_C) z_\text{prev} + \sum_{i=0}^{C-1} \exp(\Lambda_C - \Lambda_i) \phi(k_i)
\tag{4.17}
\]

where $\Lambda_C$ is the cumulative log-decay at the end of the chunk.

\subsection{Bidirectional Mode}

For encoders, we run the chunked scan twice:

\[
O^\text{fwd} = \text{Scan}(Q, K, V, \lambda) \quad \text{(left to right)}
\]

\[
O^\text{bwd} = \text{Scan}(Q, K, V, \lambda) \quad \text{(right to left)}
\]

where the backward scan flips the sequence, runs the forward scan, then 
flips back.

Final output:

\[
O = \frac{1}{2} (O^\text{fwd} + O^\text{bwd})
\tag{4.18}
\]

\subsection{Algorithm}

\textbf{Algorithm 1:} Chunked decay scan
\begin{lstlisting}
Input:  Q, K, V $\in$ R^{B$\times$H$\times$L$\times$d}, $\lambda$ $\in$ R^{B$\times$H$\times$L}
        chunk_size C = 32
Output: O $\in$ R^{B$\times$H$\times$L$\times$d}

 1: log_$\lambda$ $\leftarrow$ log(clamp($\lambda$, 1e-4, 1-1e-4))
 2:   $\leftarrow$ cumsum(log_$\lambda$, dim=-1)
 3: Pad L to multiple of C
 4: Reshape into (B, H, num_chunks, C, d)
 5: S $\leftarrow$ zeros(B, H, d, d)        Inter-chunk state
 6: z $\leftarrow$ zeros(B, H, d, 1)
 7: for c in 0..num_chunks-1:
 8:      _c $\leftarrow$  [:,:,c*C:(c+1)*C]           Cumulative in chunk
 9:     D_intra $\leftarrow$ exp( _c[:,:,:,None] -  _c[:,:,None,:])
10:     D_intra = tril(D_intra)               Causal
11:     A_intra $\leftarrow$ Q_c @ K_c^T * D_intra
12:     num_intra $\leftarrow$ A_intra @ V_c
13:     denom_intra $\leftarrow$ A_intra.sum(-1)
14:
15:     D_inter $\leftarrow$ exp( _c)    (B, H, C, 1)
16:     num_inter $\leftarrow$ Q_c @ S * D_inter
17:     denom_inter $\leftarrow$ Q_c @ z * D_inter
18:
19:     O_c $\leftarrow$ (num_intra + num_inter) / (denom_intra + denom_inter + $\epsilon$)
20:
21:     # Update state for next chunk
22:      _end $\leftarrow$  _c[:, :, -1:]    (B, H, 1)
23:     decay_to_end $\leftarrow$ exp( _end -  _c)    (B, H, C, 1)
24:     k_to_end $\leftarrow$ K_c * decay_to_end
25:     S $\leftarrow$ S * exp( _end).squeeze(-1) + k_to_end^T @ V_c
26:     z $\leftarrow$ z * exp( _end).transpose + k_to_end.sum(dim=-2)
27: return O
\end{lstlisting}

\subsection{Complexity Analysis}

For a single layer with $L$ tokens:

\begin{table}[H]
\centering
\small
\begin{tabular}{lll}
\toprule
\textbf{Component} & \textbf{Time} & \textbf{Memory} \\
\midrule
Projections & $\mathcal{O}(L d^2)$ & $\mathcal{O}(L d)$ \\
Decay gate & $\mathcal{O}(L d H)$ & $\mathcal{O}(L H)$ \\
Chunk cumsum & $\mathcal{O}(L)$ & $\mathcal{O}(L)$ \\
Intra-chunk & $\mathcal{O}(L \cdot C)$ & $\mathcal{O}(L \cdot C)$ \\
Inter-chunk & $\mathcal{O}(L \cdot d_h)$ & $\mathcal{O}(d_h^2)$ \\
State update & $\mathcal{O}(L \cdot d_h^2 / C)$ & $\mathcal{O}(d_h^2)$ \\
\textbf{Total} & $\mathcal{O}(L d^2 + L \cdot C + L \cdot d_h^2 / C)$ & $\mathcal{O}(L d + d_h^2)$ \\
\bottomrule
\end{tabular}
\end{table}

\textbf{Key observation:} The state memory is $\mathcal{O}(d_h^2)$, independent 
of $L$. In contrast, CHELA-SLAY's causal cumsum uses $\mathcal{O}(L \cdot 
2F \cdot d\_h)$.

For $d_h = 64$, $C = 32$:

\begin{table}[H]
\centering
\small
\begin{tabular}{lll}
\toprule
\textbf{Metric} & \textbf{DecayLinear} & \textbf{CHELA-SLAY causal} \\
\midrule
State memory & $64^2 = 4096$ floats = \textbf{16 KB} & $128 \times 256 \times 64$ per layer = \textbf{8 GB} \\
\bottomrule
\end{tabular}
\end{table}

\textbf{500,000$\times$ less memory} at $L = 128$.

\subsection{Parameter Analysis}

\begin{table}[H]
\centering
\small
\begin{tabular}{lll}
\toprule
\textbf{Component} & \textbf{Baseline} & \textbf{DecayLinear} \\
\midrule
Q, K, V, O & $4 d^2$ & $4 d^2$ (frozen) \\
LoRA & 0 & $8 r_L d$ \\
Decay gate $W_g$ & 0 & $d H$ \\
Decay gate $b_g$ & 0 & $H$ \\
Per-head $\ell$ & 0 & $H$ \\
\bottomrule
\end{tabular}
\end{table}

For $d = 768$, $H = 12$, $r_L = 16$:

\begin{table}[H]
\centering
\small
\begin{tabular}{ll}
\toprule
\textbf{Component} & \textbf{Params} \\
\midrule
LoRA & 98,304 \\
$W_g$ & 9,216 \\
$b_g$ & 12 \\
$\ell$ & 12 \\
\textbf{Total added} & ~108,000 \\
\bottomrule
\end{tabular}
\end{table}

\textbf{Relative to baseline attention:} +4.6\%.

\subsection{Step-0 Preservation}

\textbf{Claim.} At initialization ($W_g = 0$, $b_g = 0$, LoRA $B = 0$), the 
effective decay is \textbf{constant per head} at value $\sigma(\ell_h)$, and 
the linear attention output is a specific linear combination of V.

\textbf{This is NOT equivalent to softmax attention} -- the linear stream computes 
a different output. But the LoRA adapters can correct the discrepancy 
during recovery.

\subsection{Weight Transfer}

Base weights $W_Q, W_K, W_V, W_O$ are copied from the pretrained attention. 
The decay gate and per-head bias are initialized as above; the LoRA adapters 
are zero.

